\documentclass[10pt,a4paper]{amsart}
\usepackage{amsmath,amssymb,amsthm,graphicx}
\usepackage{verbatim}
\newenvironment{reference}{\comment}{\endcomment}
%\newenvironment{reference}{}{}
\newenvironment{slogan}{\comment}{\endcomment}
\newenvironment{history}{\comment}{\endcomment}
% For commutative diagrams we use Xy-pic
\usepackage[all]{xy}

% We use 2cell for 2-commutative diagrams.
\xyoption{2cell}
% Theorem environments.
%
\theoremstyle{plain}
\newtheorem{theorem}[subsection]{Theorem}
\newtheorem{proposition}[subsection]{Proposition}
\newtheorem{lemma}[subsection]{Lemma}

\theoremstyle{definition}
\newtheorem{definition}[subsection]{Definition}
\newtheorem{example}[subsection]{Example}
\newtheorem{exercise}[subsection]{Exercise}
\newtheorem{situation}[subsection]{Situation}

\theoremstyle{remark}
\newtheorem{remark}[subsection]{Remark}
\newtheorem{remarks}[subsection]{Remarks}

\numberwithin{equation}{subsection}

% Macros
%
\def\lim{\mathop{\mathrm{lim}}\nolimits}
\def\colim{\mathop{\mathrm{colim}}\nolimits}
\def\Spec{\mathop{\mathrm{Spec}}}
\def\Hom{\mathop{\mathrm{Hom}}\nolimits}
\def\Ext{\mathop{\mathrm{Ext}}\nolimits}
\def\SheafHom{\mathop{\mathcal{H}\!\mathit{om}}\nolimits}
\def\SheafExt{\mathop{\mathcal{E}\!\mathit{xt}}\nolimits}
\def\Sch{\mathit{Sch}}
\def\Mor{\mathop{\mathrm{Mor}}\nolimits}
\def\Ob{\mathop{\mathrm{Ob}}\nolimits}
\def\Sh{\mathop{\mathit{Sh}}\nolimits}
\def\NL{\mathop{N\!L}\nolimits}
\def\CH{\mathop{\mathrm{CH}}\nolimits}
\def\proetale{{pro\text{-}\acute{e}tale}}
\def\etale{{\acute{e}tale}}
\def\QCoh{\mathit{QCoh}}
\def\Ker{\mathop{\mathrm{Ker}}}
\def\Im{\mathop{\mathrm{Im}}}
\def\Coker{\mathop{\mathrm{Coker}}}
\def\Coim{\mathop{\mathrm{Coim}}}

% Boxtimes
%
\DeclareMathSymbol{\boxtimes}{\mathbin}{AMSa}{"02}

%
% Macros for moduli stacks/spaces
%
\def\QCohstack{\mathcal{QC}\!\mathit{oh}}
\def\Cohstack{\mathcal{C}\!\mathit{oh}}
\def\Spacesstack{\mathcal{S}\!\mathit{paces}}
\def\Quotfunctor{\mathrm{Quot}}
\def\Hilbfunctor{\mathrm{Hilb}}
\def\Curvesstack{\mathcal{C}\!\mathit{urves}}
\def\Polarizedstack{\mathcal{P}\!\mathit{olarized}}
\def\Complexesstack{\mathcal{C}\!\mathit{omplexes}}
% \Pic is the operator that assigns to X its picard group, usage \Pic(X)
% \Picardstack_{X/B} denotes the Picard stack of X over B
% \Picardfunctor_{X/B} denotes the Picard functor of X over B
\def\Pic{\mathop{\mathrm{Pic}}\nolimits}
\def\Picardstack{\mathcal{P}\!\mathit{ic}}
\def\Picardfunctor{\mathrm{Pic}}
\def\Deformationcategory{\mathcal{D}\!\mathit{ef}}
\usepackage{cleveref}
\crefname{originallemma}{original lemma}{original lemmas}
\Crefname{originallemma}{Original lemma}{Original lemmas}
\makeatletter
\newlabel{lemma-local-source-target}{{lemma-local-source-target}{source}{}{sourcepointer-1}{}}
\newlabel{lemma-local-source-target@cref}{{[originallemma][0][]lemma-local-source-target}{[1][0][]source}{}{}{}}
\makeatother
\crefname{originalreference}{original reference}{original references}
\Crefname{originalreference}{Original reference}{Original references}
\makeatletter
\newlabel{spaces-morphisms-lemma-where-smooth}{{spaces-morphisms-lemma-where-smooth}{source}{}{sourcepointer-2}{}}
\newlabel{spaces-morphisms-lemma-where-smooth@cref}{{[originalreference][0][]spaces-morphisms-lemma-where-smooth}{[1][0][]source}{}{}{}}
\makeatother
\crefname{originalreference}{original reference}{original references}
\Crefname{originalreference}{Original reference}{Original references}
\makeatletter
\newlabel{stacks-properties-lemma-open-substacks}{{stacks-properties-lemma-open-substacks}{source}{}{sourcepointer-3}{}}
\newlabel{stacks-properties-lemma-open-substacks@cref}{{[originalreference][0][]stacks-properties-lemma-open-substacks}{[1][0][]source}{}{}{}}
\makeatother
\crefname{originalreference}{original reference}{original references}
\Crefname{originalreference}{Original reference}{Original references}
\makeatletter
\newlabel{stacks-properties-lemma-topology-points}{{stacks-properties-lemma-topology-points}{source}{}{sourcepointer-4}{}}
\newlabel{stacks-properties-lemma-topology-points@cref}{{[originalreference][0][]stacks-properties-lemma-topology-points}{[1][0][]source}{}{}{}}
\makeatother
\title{The largest open substack where a morphism of algebraic stacks is smooth, and its base change behaviour}
\author{The Stacks Project Authors}
\date{stacks-project:stacks-morphisms.tex:local-snapshot-20261012; GNU FDL 1.2 or later}
\begin{document}
\maketitle
\noindent\textit{Editorial note.} All delivered mathematical text is the author's own, copied from the pinned Stacks Project chapter source sources/stacks-algebraic-stacks-morphisms/pinned-original-source.tex; no statement or proof step was rewritten, reordered or removed.\par\medskip
\noindent\textit{Editorial note.} Preamble provenance: the Stacks Project keeps one single shared preamble for every chapter; this card takes it directly from the staged sources/stacks-algebraic-stacks-morphisms/preamble.tex as three verbatim spans (lines 8-12, 14-18 and 154-217), so no derived file was written and nothing was added or redefined. The wrapper class and the base packages are supplied by the assembler.\par\medskip
\noindent\textit{Editorial note.} Bibliography provenance: neither the counted statement nor its proof contains a citation (this lane rescans the whole span before writing the card), so the card carries no bibliography and no entry was reconstructed.\par\medskip
\noindent\textit{Editorial note.} Selection provenance: this card was opened under the file-based control contract (CONTROL.md), with the mechanical pre-card gate run on its statement and proof interval before the card was written.\par\medskip
\noindent\textit{Allowed standard prerequisites (editorial).}\begin{itemize}
\item The existence of a maximal open subspace over which a morphism of algebraic spaces is smooth. Invoked inside the counted proof (source lines 7676-7678) as a cross-chapter pointer. Supplies the open subspace of the presenting space on which the induced morphism is smooth.
\item The correspondence between open substacks and open subspaces of a presentation. Invoked inside the counted proof (source lines 7679-7682) as cross-chapter pointers. Descends the open subspace to an open substack of the source.
\item The equivalence characterizing smoothness of a morphism of algebraic stacks through a presentation. Invoked inside the counted proof (source lines 7683-7685). Turns smoothness of the descended morphism into smoothness of the restriction of the given morphism, and gives maximality.
\item The statement on maximal open subschemes where a morphism of algebraic spaces is smooth. Invoked at the end of the counted proof (source lines 7691-7694) as a cross-chapter pointer. Gives clauses (2) and (3) from the corresponding statements for algebraic spaces.
\end{itemize}\medskip
\ifdefined\hypertarget\hypertarget{sourcepointer-1}{}\fi
\noindent\textit{Source reference.} Original lemma lemma-local-source-target: sources/stacks-algebraic-stacks-morphisms/pinned-original-source.tex, source lines 3237--3237. Navigation pointer only: the counted text refers to the Stacks item lemma-local-source-target, which is an existing labeled result of this same chapter and is not part of this card. The pointer carries the target's own label as its number and its definition line as its span; its content is not reproduced here.\par\medskip
\ifdefined\hypertarget\hypertarget{sourcepointer-2}{}\fi
\noindent\textit{Source reference.} Original reference spaces-morphisms-lemma-where-smooth: sources/stacks-algebraic-stacks-morphisms/pinned-original-source.tex, source lines 7679--7679. Navigation pointer only: the counted text refers to the Stacks item spaces-morphisms-lemma-where-smooth, which belongs to another chapter of the project and is not part of this card. The pointer carries the target's own label as its number because this chapter's staged snapshot does not contain the other chapter's tag table; its content is not reproduced here.\par\medskip
\ifdefined\hypertarget\hypertarget{sourcepointer-3}{}\fi
\noindent\textit{Source reference.} Original reference stacks-properties-lemma-open-substacks: sources/stacks-algebraic-stacks-morphisms/pinned-original-source.tex, source lines 7176--7176. Navigation pointer only: the counted text refers to the Stacks item stacks-properties-lemma-open-substacks, which belongs to another chapter of the project and is not part of this card. The pointer carries the target's own label as its number because this chapter's staged snapshot does not contain the other chapter's tag table; its content is not reproduced here.\par\medskip
\ifdefined\hypertarget\hypertarget{sourcepointer-4}{}\fi
\noindent\textit{Source reference.} Original reference stacks-properties-lemma-topology-points: sources/stacks-algebraic-stacks-morphisms/pinned-original-source.tex, source lines 2128--2128. Navigation pointer only: the counted text refers to the Stacks item stacks-properties-lemma-topology-points, which belongs to another chapter of the project and is not part of this card. The pointer carries the target's own label as its number because this chapter's staged snapshot does not contain the other chapter's tag table; its content is not reproduced here.\par\medskip
\begin{lemma}
\label{lemma-where-smooth}
Let $f : \mathcal{X} \to \mathcal{Y}$ be a morphism of algebraic stacks.
There is a largest open substack $\mathcal{U} \subset \mathcal{X}$
such that $f|_\mathcal{U} : \mathcal{U} \to \mathcal{Y}$ is smooth.
Moreover, formation of this open commutes with
\begin{enumerate}
\item precomposing by smooth morphisms,
\item base change by morphisms which are flat and locally of
finite presentation,
\item base change by flat morphisms provided $f$ is locally of
finite presentation.
\end{enumerate}
\end{lemma}

\begin{proof}
Choose a commutative diagram
$$
\xymatrix{
U \ar[d]_a \ar[r]_h & V \ar[d]^b \\
\mathcal{X} \ar[r]^f & \mathcal{Y}
}
$$
where $U$ and $V$ are algebraic spaces, the vertical arrows are smooth,
and $a : U \to \mathcal{X}$ surjective. There is a maximal open subspace
$U' \subset U$ such that $h_{U'} : U' \to V$ is smooth, see
Morphisms of Spaces, Lemma \ref{spaces-morphisms-lemma-where-smooth}.
Let $\mathcal{U} \subset \mathcal{X}$ be the open substack
corresponding to the image of $|U'| \to |\mathcal{X}|$
(Properties of Stacks, Lemmas \ref{stacks-properties-lemma-topology-points} and
\ref{stacks-properties-lemma-open-substacks}).
By the equivalence in Lemma \ref{lemma-local-source-target}
we find that $f|_\mathcal{U} : \mathcal{U} \to \mathcal{Y}$ is smooth
and that $\mathcal{U}$ is the largest open substack with this
property.

\medskip\noindent
Assertion (1) follows from the fact that being smooth
is smooth local on the source (this property was used to even define
smooth morphisms of algebraic stacks).
Assertions (2) and (3) follow from the case of algebraic spaces, see
Morphisms of Spaces, Lemma \ref{spaces-morphisms-lemma-where-smooth}.
\end{proof}

\end{document}
